% ============================================================
\chapter{Radon Measures and Riesz Representation}
\label{ch:radon_riesz}
% ============================================================

\section{Introduction}

This chapter lies at the interface of measure theory and topology. We study
\emph{Radon measures}\index{Radon measure} on locally compact spaces and establish the
celebrated Riesz representation theorem\index{Riesz representation theorem}, which
identifies positive linear functionals on the space of compactly supported continuous
functions with Radon measures.

\begin{intuition}[title=\textbf{Linking topology and measure}]
Radon measures are measures that ``respect'' the topology: they are determined by their
values on compact sets and open sets. The Riesz theorem says that every ``reasonable''
way of integrating compactly supported continuous functions comes from such a measure.
This is the fundamental bridge between functional analysis and measure theory.
\end{intuition}

% ------------------------------------------------------------
\section{Locally compact spaces}
% ------------------------------------------------------------

\begin{definition}[Locally compact space]
A Hausdorff topological space $X$ is \emph{locally compact}\index{locally compact}
if every point has a compact neighborhood, i.e.\ for every $x \in X$, there exist
an open set $U$ and a compact set $K$ with $x \in U \subset K$.
\end{definition}

\begin{example}
\begin{enumerate}
    \item $\R^n$ is locally compact.
    \item Every compact space is locally compact.
    \item Every open subset of a locally compact space is locally compact.
    \item An infinite-dimensional Banach space is not locally compact.
\end{enumerate}
\end{example}

\begin{definition}[Support and compactly supported functions]
For a continuous function $f : X \to \R$, the \emph{support} of $f$ is
$\mathrm{supp}(f) = \overline{\{x \in X : f(x) \neq 0\}}$. We denote by
$C_c(X)$\index{$C_c(X)$} the space of continuous functions with compact support.
\end{definition}

\begin{lemma}[Urysohn for locally compact spaces]
\label{lem:urysohn}
Let $X$ be locally compact Hausdorff, $K$ compact and $U$ open with
$K \subset U$. Then there exists $f \in C_c(X)$ with $0 \leq f \leq 1$,
$f|_K = 1$, and $\mathrm{supp}(f) \subset U$.
\end{lemma}

\begin{proposition}[Partition of unity]
\label{prop:partition_unity}
Let $X$ be locally compact Hausdorff and let $K$ be a compact set covered by open
sets $U_1, \ldots, U_n$. There exist $\varphi_1, \ldots, \varphi_n \in C_c(X)$ with
$0 \leq \varphi_i \leq 1$, $\mathrm{supp}(\varphi_i) \subset U_i$, and
$\sum_{i=1}^n \varphi_i = 1$ on $K$.
\end{proposition}

% ------------------------------------------------------------
\section{Radon measures}
% ------------------------------------------------------------

\begin{definition}[Regular Borel measure]
\label{def:regular}
Let $X$ be a locally compact Hausdorff space. A positive Borel measure $\mu$ on $X$
is called:
\begin{itemize}
    \item \emph{outer regular}\index{outer regularity} if for every
    $B \in \mathcal{B}(X)$:
    $\mu(B) = \inf\{\mu(U) : U \text{ open},\, B \subset U\}$;
    \item \emph{inner regular}\index{inner regularity} if for every open set $U$:
    $\mu(U) = \sup\{\mu(K) : K \text{ compact},\, K \subset U\}$;
    \item \emph{regular} if it is both outer and inner regular.
\end{itemize}
\end{definition}

\begin{definition}[Radon measure]
\label{def:radon}
\index{Radon measure}
A \emph{Radon measure} on a locally compact Hausdorff space $X$ is a Borel measure
$\mu$ on $X$ that is:
\begin{enumerate}
    \item[(i)] \emph{locally finite}: for every $x \in X$, there exists an open
    neighborhood $U$ of $x$ with $\mu(U) < \infty$;
    \item[(ii)] \emph{inner regular on open sets}: for every open set $U$,
    $\mu(U) = \sup\{\mu(K) : K \text{ compact},\, K \subset U\}$;
    \item[(iii)] \emph{outer regular}: for every Borel set $B$,
    $\mu(B) = \inf\{\mu(U) : U \text{ open},\, B \subset U\}$.
\end{enumerate}
\end{definition}

\begin{example}
\begin{enumerate}
    \item Lebesgue measure on $\R^n$ is a Radon measure.
    \item Every Dirac measure $\delta_a$ is a Radon measure.
    \item The counting measure on a discrete space is a Radon measure.
    \item The counting measure on $\R$ is \emph{not} a Radon measure (it is not
    locally finite).
\end{enumerate}
\end{example}

\begin{proposition}
Every finite Radon measure on a locally compact Hausdorff space is regular (inner
regular on \emph{all} Borel sets, not just open sets).
\end{proposition}

% ------------------------------------------------------------
\section{Positive linear functionals}
% ------------------------------------------------------------

\begin{definition}[Positive linear functional]
A linear functional $\Lambda : C_c(X) \to \R$ is called \emph{positive} if:
\[
    f \geq 0 \implies \Lambda(f) \geq 0.
\]
\end{definition}

\begin{remark}
Every positive linear functional on $C_c(X)$ is automatically continuous for the
topology of uniform convergence on compact sets (by positivity,
$\abs{\Lambda(f)} \leq \Lambda(\norm{f}_\infty \cdot \ind_K)$ where $K$ contains
the support of $f$).
\end{remark}

% ------------------------------------------------------------
\section{The Riesz representation theorem}
% ------------------------------------------------------------

\begin{theorem}[Riesz-Markov-Kakutani]
\label{thm:riesz}
\index{Riesz-Markov-Kakutani theorem}
Let $X$ be a locally compact Hausdorff space and let $\Lambda : C_c(X) \to \R$ be a
positive linear functional. Then there exists a unique Radon measure $\mu$ on $X$
such that:
\[
    \Lambda(f) = \int_X f\, d\mu, \qquad \forall f \in C_c(X).
\]
\end{theorem}

\begin{proof}[Sketch of proof]
The proof proceeds in several steps:

\emph{Step 1.} Define $\mu$ on open sets. For $U$ open, set:
\[
    \mu(U) = \sup\{\Lambda(f) : f \in C_c(X),\, 0 \leq f \leq 1,\,
    \mathrm{supp}(f) \subset U\}.
\]

\emph{Step 2.} Extend to Borel sets by outer regularity:
\[
    \mu(B) = \inf\{\mu(U) : U \text{ open},\, B \subset U\}.
\]

\emph{Step 3.} Verify that $\mu$ is a measure. Use Urysohn's lemma and partitions of
unity to verify $\sigma$-additivity.

\emph{Step 4.} Verify that $\Lambda(f) = \int f\, d\mu$. Approximate $f$ by simple
functions and use the regularity of $\mu$.

\emph{Step 5.} Uniqueness. If $\mu'$ is another Radon measure with
$\Lambda(f) = \int f\, d\mu'$, then $\mu$ and $\mu'$ agree on open sets (by
Urysohn's lemma), hence on all Borel sets (by regularity).
\end{proof}

\begin{attention}[title=\textbf{Conditions of the theorem}]
The space $X$ must be locally compact and Hausdorff. The functional must be defined
on $C_c(X)$ (compactly supported functions), not $C_b(X)$ (bounded functions). For
more general spaces, different versions of the theorem are needed.
\end{attention}

% ------------------------------------------------------------
\section{Support of a measure}
% ------------------------------------------------------------

\begin{definition}[Support of a measure]
\label{def:support_measure}
\index{support of a measure}
Let $\mu$ be a Radon measure on a locally compact Hausdorff space $X$. The
\emph{support} of $\mu$ is the smallest closed set $F$ such that
$\mu(X \setminus F) = 0$:
\[
    \mathrm{supp}(\mu) = X \setminus \bigcup\{U \text{ open} : \mu(U) = 0\}.
\]
\end{definition}

\begin{proposition}
Let $\mu$ be a Radon measure on $X$. Then:
\begin{enumerate}
    \item $x \in \mathrm{supp}(\mu)$ if and only if $\mu(U) > 0$ for every open
    neighborhood $U$ of $x$.
    \item $\mathrm{supp}(\mu) = \varnothing$ if and only if $\mu = 0$.
    \item If $f \in C_c(X)$ and
    $\mathrm{supp}(f) \cap \mathrm{supp}(\mu) = \varnothing$, then
    $\int f\, d\mu = 0$.
\end{enumerate}
\end{proposition}

\begin{proof}
(1) $x \notin \mathrm{supp}(\mu)$ means $x$ belongs to an open set $U$ with
$\mu(U) = 0$, i.e.\ some open neighborhood of $x$ has measure zero.

(2) $\mathrm{supp}(\mu) = \varnothing$ means $\mu(X) = 0$ by inner regularity
(every compact $K$ is contained in $X \setminus \mathrm{supp}(\mu)$, a union of
open sets of measure zero).

(3) $\mathrm{supp}(f)$ is compact and contained in the open set
$X \setminus \mathrm{supp}(\mu)$, which has measure zero.
\end{proof}

\begin{example}
\begin{enumerate}
    \item $\mathrm{supp}(\delta_a) = \{a\}$.
    \item $\mathrm{supp}(\lambda|_{[0,1]}) = [0,1]$ where $\lambda$ is Lebesgue measure.
    \item The Cantor measure has the Cantor set as its support.
\end{enumerate}
\end{example}

% ------------------------------------------------------------
\section{Riesz theorem for signed measures}
% ------------------------------------------------------------

\begin{theorem}
\label{thm:riesz_signed}
Let $X$ be locally compact Hausdorff. Every continuous linear functional
$\Lambda : C_0(X) \to \R$ (where $C_0(X)$ is the space of continuous functions
vanishing at infinity, with the uniform norm) can be uniquely written as:
\[
    \Lambda(f) = \int_X f\, d\mu
\]
where $\mu$ is a regular signed Radon measure with $\norm{\Lambda} = \abs{\mu}(X)$.
\end{theorem}

\begin{keyformulas}[title=\textbf{Riesz correspondence}]
\[
\boxed{
    \bigl(C_c(X)\bigr)^*_+ \;\longleftrightarrow\;
    \{\text{positive Radon measures on } X\}
}
\]
\[
\boxed{
    \bigl(C_0(X)\bigr)^* \;\longleftrightarrow\;
    \{\text{regular signed Radon measures on } X\}
}
\]
\end{keyformulas}

% ------------------------------------------------------------
\section{Exercises}
% ------------------------------------------------------------

\begin{exercise}
Show that Lebesgue measure on $\R$ is a Radon measure.
\emph{Hint}: verify the three conditions in Definition~\ref{def:radon}.
\end{exercise}

\begin{exercise}
Let $X = \R$ and $\Lambda(f) = f(0)$ for $f \in C_c(\R)$. Identify the Radon measure
$\mu$ such that $\Lambda(f) = \int f\, d\mu$.
\end{exercise}

\begin{exercise}
Let $\mu$ be a Radon measure on $\R^n$ and $K$ a compact set with $\mu(K) = 0$. Show
that for every $\varepsilon > 0$, there exists an open set $U \supset K$ with
$\mu(U) < \varepsilon$.
\end{exercise}

\begin{exercise}
Let $X$ be locally compact Hausdorff and $\mu$ a finite Radon measure. Show that for
every $B \in \mathcal{B}(X)$ and every $\varepsilon > 0$, there exist a compact set
$K \subset B$ and an open set $U \supset B$ with $\mu(U \setminus K) < \varepsilon$.
\end{exercise}

\begin{exercise}[Support and density]
Let $f : \R \to [0, +\infty)$ be continuous and $\mu(A) = \int_A f\, d\lambda$.
Show that $\mathrm{supp}(\mu) = \overline{\{x \in \R : f(x) > 0\}}$.
\end{exercise}

\begin{exercise}
Let $X$ be locally compact Hausdorff and let $\mu, \nu$ be two Radon measures such
that $\int f\, d\mu = \int f\, d\nu$ for every $f \in C_c(X)$. Show that $\mu = \nu$.
Is this statement true if $C_c(X)$ is replaced by a proper subspace?
\end{exercise}
